\chapter{Trade-Flow Features}\label{ch:rs:trade-flow-features}

On 6 May 2010 the Dow Jones Industrial Average made the biggest one-day point decline in its history, 998.5 points, and for a few
minutes a trillion dollars of market value vanished. Within a year three researchers argued that a measure of order-flow toxicity, \omterm{def:rs:trade-flow-features:vpin}{VPIN},
had reached historically high readings before the crash and might help signal the next turmoil. In 2014 another study found \omterm{def:rs:trade-flow-features:vpin}{VPIN} a poor
\omterm{def:rs:anatomy-of-a-predictor:predictor}{predictor} of short-run volatility. Trade flow is the market's record of who wanted to trade urgently, and at what cost; turning it into
features raises three questions this chapter answers in turn: which side initiated each trade, how the signs of trades depend
on each other, and when flow is informed. Its data are three hours of \texttt{firm.tape} with a ten-minute episode in which the
efficient price moves eight times faster and informed orders rise from 15\% to 45\% of the flow: a toxic episode whose truth the
simulator records, so that every measure can be checked against it.

\section{Signing trades}

\begin{definition}[Trade sign, tick rule, quote rule, Lee--Ready algorithm]\label{def:rs:trade-flow-features:sign}
The \emph{trade sign}\index{trade sign} is $+1$ when a trade was initiated by a buyer (an order that took the ask) and $-1$ when
initiated by a seller. The \emph{tick rule}\index{tick rule} signs a trade $+1$ if its price is above the last different trade
price and $-1$ if below. The \emph{quote rule}\index{quote rule} signs it $+1$ if its price is above the prevailing mid and $-1$
if below. The \emph{Lee--Ready algorithm}\index{Lee--Ready algorithm} applies the quote rule and falls back on the tick rule
for trades at the mid.
\end{definition}

Most public trade data do not say who initiated a trade, and the rules infer it. Lee and Ready (1991) identified the two
problems of the \omterm{def:rs:trade-flow-features:sign}{quote rule}: quotes may be recorded ahead of the trades that triggered them, and trades inside the spread are not
readily classifiable. Timing is the practical enemy. On the simulated tape, where every trade happens at the touch, the \omterm{def:rs:trade-flow-features:sign}{quote
rule} with the quotes prevailing at the \omterm{def:rs:trade-flow-features:sign}{trade signs} 100\% of trades correctly; with quotes 0.2 seconds stale, 99.1\%; one second
stale, 96.3\% (Lee--Ready 97.1\%); the \omterm{def:rs:trade-flow-features:sign}{tick rule}, which needs no quotes, 95.2\%
(\cref{fig:rs:trade-flow-features:signing}). Real data do worse, since real trades happen inside the spread, in hidden orders and
in auctions.

\begin{omfigure}
\begin{tikzpicture}[font=\small,node distance=6mm and 10mm,
  box/.style={draw=omDef,fill=omDef!8,rounded corners=2pt,align=center,inner sep=4pt,minimum height=8mm},
  res/.style={draw=omThm,fill=omThm!10,rounded corners=2pt,align=center,inner sep=4pt,minimum height=8mm},
  arr/.style={-Stealth,thick,draw=black!60}]
\node[box] (q) {quotes as of the trade\\(with the clock lag decided)};
\node[box,below=of q] (c) {price vs.\ mid};
\node[res,left=14mm of c] (buy) {buy ($+1$)};
\node[res,right=14mm of c] (sell) {sell ($-1$)};
\node[box,below=of c] (t) {at the mid: price vs.\ last\\different trade price};
\node[res,left=14mm of t] (tb) {uptick: buy};
\node[res,right=14mm of t] (ts) {downtick: sell};
\draw[arr] (q) -- (c);
\draw[arr] (c) -- node[above,font=\footnotesize]{above} (buy);
\draw[arr] (c) -- node[above,font=\footnotesize]{below} (sell);
\draw[arr] (c) -- node[right,font=\footnotesize]{equal} (t);
\draw[arr] (t) -- (tb);
\draw[arr] (t) -- (ts);
\end{tikzpicture}

\omcaption{The \omterm{def:rs:trade-flow-features:sign}{Lee--Ready algorithm}: the \omterm{def:rs:trade-flow-features:sign}{quote rule} first, the \omterm{def:rs:trade-flow-features:sign}{tick rule} for trades at the mid. The first box carries most of the error: on
the simulated tape a one-second clock lag costs the \omterm{def:rs:trade-flow-features:sign}{quote rule} 3.7 points of accuracy.}\label{fig:rs:trade-flow-features:leeready}
\end{omfigure}

\begin{definition}[Bulk volume classification]\label{def:rs:trade-flow-features:bvc}
\emph{Bulk volume classification}\index{bulk volume classification} (BVC) splits a \omterm{def:rs:market-data-for-research:bar}{bar}'s volume into buys and sells without
signing individual trades: the buy share is $\Phi(\Delta P/\sigma_{\Delta P})$, with $\Delta P$ the \omterm{def:rs:market-data-for-research:bar}{bar}'s price change and
$\sigma_{\Delta P}$ its standard deviation across \omterm{def:rs:market-data-for-research:bar}{bars}.
\end{definition}

BVC was introduced with \omterm{def:rs:trade-flow-features:vpin}{VPIN} for data too coarse to sign trade by trade. It pays for that: on one-minute \omterm{def:rs:market-data-for-research:bar}{bars} of the simulated tape
it classifies 80.4\% of the volume correctly, against 95\% or more for the trade-level rules, and Chakrabarty, Pascual and Shkilko
found on a large sample of US stocks that the \omterm{def:rs:trade-flow-features:sign}{tick rule} and Lee--Ready classify significantly better than BVC.

\begin{omfigure}
\begin{tikzpicture}
\begin{axis}[width=11cm,height=5.4cm,xbar,bar width=8pt,xlabel={share classified correctly},tick label style={font=\small},
  label style={font=\small},xmin=0.75,xmax=1.01,ymin=-0.6,ymax=6.6,ytick={0,1,2,3,4,5,6},
  yticklabels={{quote, on time},{quote, 0.2 s late},{Lee--Ready, 0.2 s late},{quote, 1 s late},{Lee--Ready, 1 s late},{tick rule},{BVC, 1-minute bars}},
  y dir=reverse,grid=major,grid style={gray!25},table/col sep=comma]
\addplot[fill=omDef!55,draw=omDef] table[x=accuracy,y=k]{figdata/research/09-trade-flow-features/signing.csv};
\end{axis}
\end{tikzpicture}

\omcaption{Share of the simulated tape's executions signed correctly by each rule (by volume for BVC), against the simulator's
true signs. Stale quotes cost the \omterm{def:rs:trade-flow-features:sign}{quote rule} 3.7 points at one second; BVC on one-minute \omterm{def:rs:market-data-for-research:bar}{bars} gets 80\% of the volume right. Data:
\texttt{firm.tape}, seed 9.}\label{fig:rs:trade-flow-features:signing}
\end{omfigure}

\section{Order-sign autocorrelation}

\begin{definition}[Order-sign autocorrelation, metaorder]\label{def:rs:trade-flow-features:acf}
The \emph{order-sign autocorrelation}\index{order-sign autocorrelation} is the autocorrelation function of the signs of successive
aggressive orders. A \emph{metaorder}\index{metaorder} is a large order executed as a sequence of smaller child orders over time.
\end{definition}

Order signs are not independent: buys follow buys. Lillo and Farmer (2004) found on the London Stock Exchange that the
autocorrelation of order signs decays roughly as a power of the lag with exponent 0.6, a Hurst exponent of 0.7 (Book 4,
chapter 17): long memory. The explanation they and Mike later gave is the splitting of \omterm{def:rs:trade-flow-features:acf}{metaorders}.

\begin{proposition}[Metaorders make order signs long-memory]\label{prop:rs:trade-flow-features:lmf}
If \omterm{def:rs:trade-flow-features:acf}{metaorders} are split into child orders of constant size, all of the \omterm{def:rs:trade-flow-features:acf}{metaorder}'s sign, and the distribution of \omterm{def:rs:trade-flow-features:acf}{metaorder} sizes has a
power-law tail $\P(V > v) \propto v^{-\alpha}$ with $1 < \alpha < 2$, the autocorrelation of child-order signs decays asymptotically as
$\tau^{-(\alpha - 1)}$, and the Hurst exponent of the sign series is $(3 - \alpha)/2$.
\end{proposition}

\admitted
\noindent The result is Lillo, Mike and Farmer's (2005); the Hurst exponent follows from the decay exponent $\gamma = \alpha - 1$ by
$H = 1 - \gamma/2$.

The simulator's noise traders draw \omterm{def:rs:trade-flow-features:acf}{metaorder} lengths from a Pareto law with tail exponent 1.5, so the proposition predicts $H = 0.75$. The
7\,273 aggressive orders of the simulated tape have sign autocorrelations of 0.20 at lag 1, 0.087 at lag 10 and 0.028 at lag 100
(\cref{fig:rs:trade-flow-features:acf}), and a Hurst exponent by aggregated variance of 0.715. Sign memory is the reason signed volume is
a \omterm{def:rs:anatomy-of-a-predictor:predictor}{predictor} of the next signs, and not only a summary of the past: a buying \omterm{def:rs:trade-flow-features:acf}{metaorder} that has shown itself will probably keep buying.

How much of that memory is a forecast of prices? The signed-flow imbalance of the last $N$ aggressive orders (the first card below) was
measured against three targets on the simulated tape: the next order's sign, the mid-price change over the next ten seconds, and, as a
control, the change over the previous ten seconds.

\begin{center}
\small
\begin{tabular}{@{}rcccc@{}}
\toprule
last $N$ orders & corr.\ with next sign & next sign right & corr.\ with next 10 s & corr.\ with last 10 s\\
\midrule
5 & 0.25 & 56.1\% & 0.16 & 0.22\\
20 & 0.24 & 58.4\% & 0.14 & 0.31\\
100 & 0.18 & 57.1\% & 0.12 & 0.17\\
\bottomrule
\end{tabular}
\end{center}

The flow forecasts the next sign well and the next price change less well, and for twenty orders it is twice as correlated with the
change that has already happened as with the one to come: much of what the flow knows, the price already shows. Order signs are
predictable without prices being so, which is Lillo and Farmer's point: the liquidity on the other side adjusts.

\begin{omfigure}
\begin{tikzpicture}
\begin{loglogaxis}[width=9.5cm,height=5.2cm,xlabel={lag (orders)},ylabel={autocorrelation of order signs},tick label style={font=\small},
  label style={font=\small},xmin=0.8,xmax=120,ymin=0.01,ymax=0.4,grid=major,grid style={gray!25},legend columns=2,
  legend style={font=\footnotesize,at={(0.5,-0.32)},anchor=north,draw=none,fill=none,
  /tikz/every even column/.append style={column sep=8pt}},table/col sep=comma]
\addplot[omDef,very thick,mark=*] table[x=lag,y=acf]{figdata/research/09-trade-flow-features/signacf.csv};
\addplot[omThm,thick,dashed,domain=1:100,samples=2] {0.2*x^(-0.5)};
\legend{simulated tape,$0.2\,\tau^{-0.5}$ (the proposition's slope)}
\end{loglogaxis}
\end{tikzpicture}

\omcaption{Autocorrelation of the signs of successive aggressive orders on the simulated tape, on logarithmic axes, with a power law of
exponent $0.5 = \alpha - 1$ for the simulator's \omterm{def:rs:trade-flow-features:acf}{metaorder} tail $\alpha = 1.5$. Data: \texttt{firm.tape}, seed 9.}\label{fig:rs:trade-flow-features:acf}
\end{omfigure}

\section{Flow toxicity}

\begin{definition}[Flow toxicity, VPIN]\label{def:rs:trade-flow-features:vpin}
\emph{Flow toxicity}\index{flow toxicity} is the adverse selection that order flow imposes on the liquidity providers who trade against it:
its aggressors are informed, and the resting side loses on average. \emph{VPIN}\index{VPIN} (volume-synchronised probability of informed
trading) cuts the trade sequence into buckets of equal volume $V$ and, at the end of each bucket, averages $|V^B_\tau - V^S_\tau|/V$ over
the last $n$ buckets, with $V^B_\tau$ and $V^S_\tau$ the bucket's buy and sell volume.
\end{definition}

\omterm{def:rs:trade-flow-features:vpin}{VPIN} measures one-sided flow. Its authors define toxicity as adverse selection of the market makers (Easley, López de Prado and O'Hara,
2012) and argue that one-sided volume is how it shows. Toxicity itself, though, is directly measurable after the fact: the mark-out of
Book~2, chapter 15, of trades against their resting side. The simulated episode separates the two cleanly. During it, the aggressors'
mark-out ten seconds after the trade goes from $-0.36$ ticks (they pay the half-spread) to $+0.96$ (they win): the flow is toxic by
definition. \omterm{def:rs:trade-flow-features:vpin}{VPIN}, over buckets of 2\,649 shares and a window of twenty, rises from 0.31 to 0.35, 13\%; its maximum over the three hours
falls outside the episode, 20 minutes after it ends (\cref{fig:rs:trade-flow-features:vpin}). The informed traders of the simulator buy when
the efficient price is above the mid and sell when below, so their flow is not one-sided over a bucket; the noise traders' \omterm{def:rs:trade-flow-features:acf}{metaorders}
are. One-sided volume and informed volume are different things.

\begin{omfigure}
\begin{tikzpicture}
\begin{axis}[name=top,scale only axis,width=9.6cm,height=3.2cm,ylabel={VPIN},tick label style={font=\small},label style={font=\small},xmin=0,xmax=180,
  ymin=0.15,ymax=0.6,grid=major,grid style={gray!25},xticklabels={},table/col sep=comma]
\fill[omThm!12] (axis cs:90,0.15) rectangle (axis cs:100,0.6);
\addplot[omDef,thick] table[x=minute,y=vpin]{figdata/research/09-trade-flow-features/vpin.csv};
\end{axis}
\begin{axis}[at={(top.south west)},anchor=north west,yshift=-3mm,scale only axis,width=9.6cm,height=3.2cm,xlabel={minutes},
  ylabel={mark-out (ticks)},tick label style={font=\small},label style={font=\small},xmin=0,xmax=180,ymin=-1.5,ymax=3.2,grid=major,
  grid style={gray!25},table/col sep=comma]
\fill[omThm!12] (axis cs:90,-1.5) rectangle (axis cs:100,3.2);
\addplot[omThm,thick] table[x=minute,y=markout]{figdata/research/09-trade-flow-features/markout.csv};
\end{axis}
\end{tikzpicture}

\omcaption{The toxic episode (shaded, minutes 90 to 100) seen by \omterm{def:rs:trade-flow-features:vpin}{VPIN} (top: 400 buckets of 2\,649 shares, window of 20) and by the
aggressors' mark-out ten seconds after each trade (bottom: rolling one-minute mean, plotted when known). The mark-out turns strongly
positive; \omterm{def:rs:trade-flow-features:vpin}{VPIN} spikes as the episode starts, but its mean over the episode rises only 13\% and its maximum falls outside it. Data: \texttt{firm.tape}, seed 9.}\label{fig:rs:trade-flow-features:vpin}
\end{omfigure}

The flash-crash debate had the same shape. Easley, López de Prado and O'Hara (2011) found historically high \omterm{def:rs:trade-flow-features:vpin}{VPIN} readings before the
crash; Andersen and Bondarenko (2014) found \omterm{def:rs:trade-flow-features:vpin}{VPIN} a poor \omterm{def:rs:anatomy-of-a-predictor:predictor}{predictor} of short-run volatility. A simulator cannot settle a debate about real
markets, but it can show a measure failing on a toxicity it was designed for, which is reason enough to validate it on data where the
truth is known before trusting it where it is not. The mark-out measure has its own cost: it needs the future (ten seconds here), so it
detects late; in the episode it first crosses its pre-episode 95th percentile 40 seconds after the start, \omterm{def:rs:trade-flow-features:vpin}{VPIN} after 16 seconds, but \omterm{def:rs:trade-flow-features:vpin}{VPIN}
also spends 6.0\% of the rest of the session above its threshold.

\section{Fitted excitation intensities}

Aggressive orders cluster in time, and a Hawkes process (Book 4, chapter 7) is the standard description: each order raises the intensity
of the next. Its branching ratio is the share of orders triggered by earlier ones. On a flat version of the simulated tape (no activity
process, no news, no informed traders) the maximum-likelihood fit recovers 0.403, against the 0.4 the simulator plants. On the chapter's
tape the fit gives 0.65. Nothing in the order-generating mechanism changed; what changed is that activity rises and falls for other
reasons, and a Hawkes process attributes every cluster to self-excitation. A fitted branching ratio is an upper bound on endogeneity
unless the baseline is allowed to vary. As features, the fitted buy and sell intensities and their difference summarise how excited each
side of the flow is.

\section{Large trades and metaorders}

\begin{definition}[Kyle's lambda]\label{def:rs:trade-flow-features:lambda}
\emph{Kyle's lambda}\index{Kyle's lambda} is the price impact per unit of signed order flow: the slope of price changes on net signed
volume over the same intervals, after the parameter of Kyle's (1985) model of an informed trader and market makers.
\end{definition}

On the simulated tape, over ten-second intervals, lambda is 0.10 ticks per 100 shares of net buying, with an $R^2$ of 14\%: far below the
53\% of \omterm{def:rs:order-book-features:ofi}{order-flow imbalance} (chapter~8), because signed trade volume misses the cancellations and additions that move the quotes, as Cont,
Kukanov and Stoikov noted for real data. Lambda is a liquidity feature (a stock whose price moves more per share traded is less liquid)
and a cost input (chapter~27). \omterm{def:rs:trade-flow-features:acf}{Metaorder} detection starts from the same flows: a run of same-sign orders of similar size at regular
intervals, an unusual share of one side over minutes, a VPIN-like imbalance on a volume clock. The chapter's memory result says why it
works: children of one parent share its sign.

\section{Predictor cards}

\begin{predictorcard}{Signed-flow imbalance}\label{pred:rs:trade-flow-features:signed}
\sfield{Definition} $\sum b_iq_i/\sum q_i$ over the aggressive orders of the last $\Delta$ seconds (or last $N$ orders), $b_i$ the sign.
\sfield{Inputs and timestamps} Trades with receive times; signs from the feed's aggressor flag or from Lee--Ready with quotes as of the
trade's receive time.
\sfield{Rationale} Order signs have long memory (\cref{prop:rs:trade-flow-features:lmf}); \omterm{def:rs:trade-flow-features:acf}{metaorders} continue.
\sfield{Horizon and half-life} The next orders; sign autocorrelation 0.20 at one order, 0.087 at ten on the simulated tape; over twenty orders it gets the next sign right 58.4\% of the time.
\sfield{Normalisation} By total volume; across instruments by average trade size.
\sfield{Failure modes} Misclassification with stale quotes; the price moves before the flow is seen (the impact is already in the price).
\sfield{Sources} Lillo and Farmer (2004); Lillo, Mike and Farmer (2005); \texttt{rs\_tradeflow.memory}.
\end{predictorcard}

\begin{predictorcard}{VPIN}\label{pred:rs:trade-flow-features:vpin}
\sfield{Definition} Mean absolute buy--sell imbalance over the last $n$ volume buckets, divided by the bucket size.
\sfield{Inputs and timestamps} Trades, signed (or BVC on \omterm{def:rs:market-data-for-research:bar}{bars}), cut into equal-volume buckets.
\sfield{Rationale} Toxic flow is one-sided (Easley, López de Prado and O'Hara, 2012).
\sfield{Horizon and half-life} Minutes to hours. On the simulated episode it rises 13\% while the aggressors' mark-out turns from $-0.36$
to $+0.96$ ticks.
\sfield{Normalisation} Bounded in $[0, 1]$; its level depends on the bucket size and window.
\sfield{Failure modes} Informed flow that is not one-sided; one-sided flow that is not informed (\omterm{def:rs:trade-flow-features:acf}{metaorders}, index rebalancing); poor
volatility prediction (Andersen and Bondarenko, 2014).
\sfield{Sources} As cited; \texttt{rs\_tradeflow.toxicity\_race}.
\end{predictorcard}

\section{Tutorial: signing and weighing the flow}\label{tut:rs:trade-flow-features:flow}

\begin{tutorial}
\textbf{Goal.} Sign the simulated tape's trades by every rule, measure the memory of order signs and fit a Hawkes process, and race \omterm{def:rs:trade-flow-features:vpin}{VPIN}
against the mark-out through the toxic episode. \textbf{End state:}
\cref{fig:rs:trade-flow-features:signing,fig:rs:trade-flow-features:acf,fig:rs:trade-flow-features:vpin}; the Hawkes branching ratios 0.65
and 0.403.
\begin{tutsteps}
\item \textbf{The signing rules.} The \omterm{def:rs:trade-flow-features:sign}{tick rule} carries the last price change forward; Lee--Ready uses the \omterm{def:rs:trade-flow-features:sign}{quote rule} and the \omterm{def:rs:trade-flow-features:sign}{tick rule} for
trades at the mid.
\omcode{code/firm/tradeflow/firm_tradeflow.py}{27}{49}{Tick rule, quote rule and Lee--Ready.}\label{lst:rs:trade-flow-features:rules}
\item \textbf{\omterm{def:rs:trade-flow-features:vpin}{VPIN}.} Trades are poured into buckets of equal volume, split where a trade straddles two buckets.
\omcode{code/firm/tradeflow/firm_tradeflow.py}{79}{101}{VPIN on volume buckets.}\label{lst:rs:trade-flow-features:vpincode}
\item \textbf{Run} \texttt{signing()}, \texttt{memory()}, \texttt{signed\_flow()}, \texttt{hawkes\_fits()}, \texttt{toxicity\_race()}, \texttt{kyle()} and
\texttt{fig\_tradeflow.py}.
\end{tutsteps}
\textbf{What to change next.} Make the informed traders one-sided (buy only while the episode lasts) and see \omterm{def:rs:trade-flow-features:vpin}{VPIN} respond; fit a Hawkes
process with a baseline that follows the activity level and watch the branching ratio fall back.
\end{tutorial}

\section{Build: the trade-flow toolkit}\label{bld:rs:trade-flow-features:tradeflow}

\begin{build}
\bfield{Purpose} Signing, memory, toxicity and impact measures on the miniature firm's trade data, for features (Books 8, 11, 12) and for
execution research (Book~10).
\bfield{Interface} \texttt{tick\_rule}, \texttt{quote\_rule}, \texttt{lee\_ready}, \texttt{bvc}, \texttt{sign\_acf}, \texttt{hurst\_aggvar},
\texttt{vpin(buy\_volume, volume, bucket, n)}, \texttt{kyle\_lambda}, \texttt{markout\_toxicity}, \texttt{aggregate\_orders}.
\bfield{Rules} Quotes passed in are those prevailing at each trade's \omterm{def:rs:point-in-time-data-and-the-biases:bitemporal}{knowledge time}; a trade that straddles a volume bucket is split
proportionally; executions of one aggressive order are aggregated before sign statistics.
\bfield{Acceptance tests} \texttt{code/firm/tradeflow/tests/}: the three rules on a hand-built sequence and BVC's symmetry; \omterm{def:rs:trade-flow-features:vpin}{VPIN} of a
balanced and of a one-sided sequence; $H = 0.5$ for independent data and a sign autocorrelation that vanishes beyond a run's length; a
known lambda recovered; order aggregation.
\bfield{Stretch} Hawkes fits with a time-varying baseline; \omterm{def:rs:trade-flow-features:acf}{metaorder} detection by change points in signed flow; mark-outs by
counterparty type.
\end{build}

\begin{omsources}
\item C.~M.~C.~Lee and M.~J.~Ready, ``Inferring trade direction from intraday data'', \emph{Journal of Finance} 46(2), 1991.
\item B.~Chakrabarty, R.~Pascual and A.~Shkilko, ``Evaluating trade classification algorithms: bulk volume classification versus the tick rule
and the Lee--Ready algorithm'', \emph{Journal of Financial Markets} 25, 2015.
\item F.~Lillo and J.~D.~Farmer, ``The long memory of the efficient market'', \emph{Studies in Nonlinear Dynamics and Econometrics} 8(3), 2004;
F.~Lillo, S.~Mike and J.~D.~Farmer, ``Theory for long memory in supply and demand'', \emph{Physical Review E} 71, 2005.
\item D.~Easley, M.~López de Prado and M.~O'Hara, ``The microstructure of the `flash crash''', \emph{Journal of Portfolio Management}
37(2), 2011; ``Flow toxicity and liquidity in a high-frequency world'', \emph{Review of Financial Studies} 25(5), 2012.
\item T.~G.~Andersen and O.~Bondarenko, ``VPIN and the flash crash'', \emph{Journal of Financial Markets} 17, 2014.
\item A.~S.~Kyle, ``Continuous auctions and insider trading'', \emph{Econometrica} 53(6), 1985.
\end{omsources}

\section{Exercises}

\begin{exercise}[$\star$]\label{exo:rs:trade-flow-features:1}
Trades at 20.00, 20.01, 20.01, 20.00, 19.99 with the quotes 20.00--20.01 throughout. Sign each by the \omterm{def:rs:trade-flow-features:sign}{tick rule} and by Lee--Ready.
\end{exercise}

\begin{exercise}[$\star$]\label{exo:rs:trade-flow-features:2}
A one-minute \omterm{def:rs:market-data-for-research:bar}{bar}'s price change is $+0.8$ standard deviations. What share of its volume does BVC classify as buys?
\end{exercise}

\begin{exercise}[$\star$]\label{exo:rs:trade-flow-features:3}
\omterm{def:rs:trade-flow-features:acf}{Metaorders} with a size tail exponent of 1.4: what decay exponent and Hurst exponent does \cref{prop:rs:trade-flow-features:lmf} predict?
\end{exercise}

\begin{exercise}[$\star\star$]\label{exo:rs:trade-flow-features:4}
Five volume buckets of 1\,000 shares have buy volumes 700, 400, 900, 500 and 200. What is \omterm{def:rs:trade-flow-features:vpin}{VPIN} over the last five buckets?
\end{exercise}

\begin{exercise}[$\star\star$]\label{exo:rs:trade-flow-features:5}
\omterm{def:rs:trade-flow-features:lambda}{Kyle's lambda} is 0.10 ticks per 100 shares. What price move does a net purchase of 3\,000 shares over ten seconds imply, and why is the
estimate noisier than one based on OFI?
\end{exercise}

\begin{exercise}[$\star\star$]\label{exo:rs:trade-flow-features:6}
Why does a Hawkes fit on data whose baseline activity varies overstate the branching ratio? Give the intuition with two regimes of
different constant intensity.
\end{exercise}

\begin{exercise}[$\star\star\star$]\label{exo:rs:trade-flow-features:7}
\emph{Coding.} With the chapter's tape, compute the share of trades signed correctly by the \omterm{def:rs:trade-flow-features:sign}{quote rule} with quotes 0.05, 0.5 and 2 seconds
late. How fast does accuracy fall with staleness?
\end{exercise}

\begin{exercise}[$\star\star\star$]\label{exo:rs:trade-flow-features:8}
\emph{Find the flaw.} ``\omterm{def:rs:trade-flow-features:vpin}{VPIN} reached its highest level of the year an hour before the crash, so \omterm{def:rs:trade-flow-features:vpin}{VPIN} predicts crashes.''
\end{exercise}

\section{Problem: Did VPIN See It Coming?}

\begin{problem}[{Weekend problem --- a toxicity measure against the truth}]\label{pb:rs:trade-flow-features:1}
The chapter's tape: three hours, a ten-minute episode from minute 90 in which the efficient price moves eight times faster and noise
trading does not rise, seed 9.

\medskip\noindent\textbf{Part I --- The truth.}
\begin{enumerate}
\item What share of the flow is informed before and during the episode?
\item What is the aggressors' ten-second mark-out, in ticks, before and during?
\item Why is the mark-out negative before the episode?
\item Is the episode's flow toxic by the definition of the chapter?
\end{enumerate}

\medskip\noindent\textbf{Part II --- \omterm{def:rs:trade-flow-features:vpin}{VPIN}.}
\begin{enumerate}[resume]
\item What bucket size and window does the chapter use, and how many buckets does the session have?
\item What are \omterm{def:rs:trade-flow-features:vpin}{VPIN}'s mean before and during the episode?
\item When does \omterm{def:rs:trade-flow-features:vpin}{VPIN} reach its maximum?
\item When does it first cross its pre-episode 95th percentile after the episode starts, and what share of the rest of the session is above
that threshold?
\item Why does \omterm{def:rs:trade-flow-features:vpin}{VPIN} barely move?
\end{enumerate}

\medskip\noindent\textbf{Part III --- The mark-out measure.}
\begin{enumerate}[resume]
\item What are its mean before and during, and when does it peak?
\item When does it first cross its threshold, and why is it later than \omterm{def:rs:trade-flow-features:vpin}{VPIN}?
\item What share of the rest of the session is above its threshold?
\item Which measure would you use to protect a market maker, and which to study toxicity after the fact?
\end{enumerate}

\medskip\noindent\textbf{Part IV --- Beyond.}
\begin{enumerate}[resume]
\item What would make informed flow one-sided, and \omterm{def:rs:trade-flow-features:vpin}{VPIN} useful?
\item What does the Hawkes fit give on this tape and on the flat tape?
\item What are \omterm{def:rs:trade-flow-features:lambda}{Kyle's lambda} and its $R^2$ at ten seconds?
\item Which of the chapter's features could detect the episode fastest, and at what cost in false alarms?
\item State the \emph{named result}: \omterm{def:rs:trade-flow-features:vpin}{VPIN}'s rise during an episode in which the informed share triples and the aggressors' mark-out turns
from $-0.36$ to $+0.96$ ticks.
\item What does the result say about validating a toxicity measure?
\item In one sentence: what is toxic flow?
\end{enumerate}
\end{problem}

\section{Interview questions}

\begin{interviewq}[$\star$ \iqroles{researcher, trader}]\label{iq:rs:trade-flow-features:1}
How would you decide whether a trade was a buy or a sell if the data do not say?
\end{interviewq}

\begin{interviewq}[$\star\star$ \iqroles{researcher}]\label{iq:rs:trade-flow-features:2}
Why are the signs of market orders autocorrelated, and why does that not make prices predictable in the same way?
\end{interviewq}

\begin{interviewq}[$\star\star$ \iqroles{trader, researcher}]\label{iq:rs:trade-flow-features:3}
What is \omterm{def:rs:trade-flow-features:vpin}{flow toxicity}, and how would you measure it for your own market-making book?
\end{interviewq}

\begin{interviewq}[$\star\star$ \iqroles{researcher, mle}]\label{iq:rs:trade-flow-features:4}
You fit a Hawkes process to trade times and get a branching ratio of 0.9. What else could explain it?
\end{interviewq}

\begin{interviewq}[$\star\star$ \iqroles{researcher}]\label{iq:rs:trade-flow-features:5}
Explain \omterm{def:rs:trade-flow-features:vpin}{VPIN} and its main criticism.
\end{interviewq}

\begin{interviewq}[$\star\star\star$ \iqroles{researcher}]\label{iq:rs:trade-flow-features:6}
Derive the relation between the tail exponent of \omterm{def:rs:trade-flow-features:acf}{metaorder} sizes and the Hurst exponent of order signs, and test it on a series.
\end{interviewq}
